\begin{ouverture}[courbe]{2}{Puissances, racines et notation scientifique}{Partie 1 · Nombres et calculs}
{Exposants · Préfixes · Ordres de grandeur · Équation $x^2=a$}
\objectif{Utiliser les puissances d’exposant entier.}
\objectif{Écrire et calculer en notation scientifique.}
\objectif{Résoudre une équation $x^2=a$ et un problème avec une racine carrée.}
\prerequis{Produits, quotients et carrés des entiers de 0 à 12.}
\end{ouverture}
\pagestyle{otMain}\essentiel\fontsize{11}{15}\selectfont
\essSection{1}{Puissances et notation scientifique}
\begin{definition}Pour $n$ entier positif, $a^n$ est le produit de $n$ facteurs égaux à $a$. Pour $a\ne0$, on pose $a^0=1$ et $a^{-n}=1/a^n$ : en divisant $10^3$ successivement par 10, on obtient $10^2$, $10^1$, $10^0=1$, puis $10^{-1}=0{,}1$.\end{definition}
\begin{propriete}Pour $a$ et $b$ non nuls et $m$, $n$ entiers relatifs :
\[a^m\times a^n=a^{m+n},\qquad \frac{a^m}{a^n}=a^{m-n},\qquad a^n\times b^n=(ab)^n.\]\end{propriete}
\begin{definition}La notation scientifique d’un nombre positif est $c\times10^p$, avec $1\le c<10$ et $p$ entier relatif. Par exemple $0{,}000056=5{,}6\times10^{-5}$.\end{definition}
\begin{savoirfaire}[Calculer en notation scientifique]Regrouper les coefficients, puis les puissances de 10. $(3\times10^5)(4\times10^{-2})=12\times10^3=1{,}2\times10^4$.\end{savoirfaire}
\begin{propriete}[Préfixes]nano : $10^{-9}$ ; micro (µ) : $10^{-6}$ ; milli : $10^{-3}$ ; kilo : $10^{3}$ ; méga : $10^{6}$ ; giga : $10^{9}$. Ainsi $1\text{ nm}=10^{-9}\text{ m}$ et $1\text{ Go}=10^9$ octets.\end{propriete}
\begin{attention}$a^m+a^n$ ne se transforme pas en $a^{m+n}$. Les règles sur les exposants concernent des produits et des quotients.\end{attention}
\clearpage
\essSection{2}{Racine carrée et équation $x^2=a$}
\begin{definition}Pour $a\ge0$, $\sqrt a$ est l’unique nombre positif ou nul dont le carré vaut $a$. Ainsi $(\sqrt a)^2=a$. La racine carrée d’un nombre négatif n’est pas définie.\end{definition}
\begin{propriete}[Résoudre $x^2=a$]Si $a>0$, l’équation a deux solutions, $-\sqrt a$ et $\sqrt a$ : l’ensemble des solutions est $\{-\sqrt a\,;\sqrt a\}$. Si $a=0$, la seule solution est 0. Si $a<0$, il n’y a aucune solution : l’ensemble des solutions est vide, noté $\varnothing$.\end{propriete}
\begin{exemple}Pour $x^2=9$, on écrit $x^2-9=(x-3)(x+3)=0$ : les solutions sont $-3$ et 3. Graphiquement, la droite d’équation $y=9$ coupe la courbe de la fonction carré en deux points d’abscisses $-3$ et 3.\end{exemple}
\begin{savoirfaire}[Arrondir une mesure]Un carré d’aire 45 cm$^2$ a pour côté $\sqrt{45}$ cm, environ $6{,}71$ cm : une longueur est positive, on ne garde pas $-\sqrt{45}$. Conserver les valeurs exactes pendant les calculs et arrondir à la fin, à la précision demandée.\end{savoirfaire}
\begin{attention}$\sqrt{(-4)^2}=4$ : une racine carrée est toujours positive ou nulle. Et $\sqrt{a+b}$ n’est pas, en général, égal à $\sqrt a+\sqrt b$.\end{attention}
\clearpage\rubrique{Découverte}
\begin{blocLivre}{ex-2-1}\begin{exercice}[1]{Comprendre un exposant}
Écrire $7\times7\times7\times7$ comme une puissance. Calculer $2^{-3}$ et $9^0$.
\end{exercice}\end{blocLivre}
\begin{blocLivre}{ex-2-2}\begin{exercice}[1]{Parenthèses et signe}
Calculer $(-5)^2$, $-5^2$, $(-2)^3$ et $(-2)^4$.
\end{exercice}\end{blocLivre}
\begin{blocLivre}{ex-2-3}\begin{exercice}[1]{Règles de puissances}
Écrire sous la forme d’une seule puissance : $3^4\times3^2$, $\dfrac{5^7}{5^3}$, $2^3\times5^3$ et $\dfrac{10^{-2}}{10^{4}}$.
\end{exercice}\end{blocLivre}
\begin{blocLivre}{ex-2-4}\begin{exercice}[1]{Écriture scientifique}
Écrire $74000000$ et $0{,}000092$ en notation scientifique.
\end{exercice}\end{blocLivre}
\begin{blocLivre}{ex-2-5}\begin{exercice}[1]{Quelques racines}
Calculer $\sqrt{144}$, $\sqrt{0{,}81}$ et $\sqrt{49/64}$.
\end{exercice}\end{blocLivre}

\clearpage\rubrique{Entraînement}
\begin{blocLivre}{ex-2-6}\begin{exercice}[2]{Produit scientifique}
Calculer $(6\times10^8)(4\times10^{-5})$ en notation scientifique.
\end{exercice}\end{blocLivre}
\begin{blocLivre}{ex-2-7}\begin{exercice}[2]{Quotient scientifique}
Calculer $\dfrac{8{,}4\times10^{-3}}{2{,}1\times10^5}$.
\end{exercice}\end{blocLivre}
\begin{blocLivre}{ex-2-8}\begin{exercice}[2]{Des préfixes}
Écrire en notation scientifique, dans l’unité indiquée : $4{,}7$ Go en octets ; 250 nm en mètres ; $0{,}6$ µs en secondes.
\end{exercice}\end{blocLivre}
\begin{blocLivre}{ex-2-9}\begin{exercice}[2]{Une équation $x^2=a$}
Résoudre $x^2=49$, $x^2=7$, $x^2=0$ et $x^2=-4$. Donner chaque ensemble de solutions.
\end{exercice}\end{blocLivre}
\begin{blocLivre}{ex-2-10}\begin{exercice}[2]{Réfuter une règle}
Réfuter $\sqrt{a+b}=\sqrt a+\sqrt b$ en choisissant deux nombres positifs simples.
\end{exercice}\end{blocLivre}

\clearpage\rubrique{Problèmes}
\begin{blocLivre}{ex-2-11}\begin{exercice}[3]{Une population qui double}
Une culture contient 800 cellules. Leur nombre double toutes les 20 minutes, selon un modèle valable pendant l’expérience. Combien de cellules après 2 heures ? Écrire le résultat en notation scientifique.
\end{exercice}\end{blocLivre}
\begin{blocLivre}{ex-2-12}\begin{exercice}[3]{Une masse minuscule}
Un lot de $4\times10^{12}$ particules a une masse totale de $1{,}8$ g. Dans ce modèle, toutes ont la même masse. Calculer la masse d’une particule en grammes, puis en kilogrammes.
\end{exercice}\end{blocLivre}
\begin{blocLivre}{ex-2-13}\begin{exercice}[3]{Une distance astronomique}
Un signal parcourt $3\times10^5$ km par seconde. Il voyage pendant 8 min 20 s. Quelle distance parcourt-il ? Donner aussi cette distance en mètres.
\end{exercice}\end{blocLivre}
\begin{blocLivre}{ex-2-14}\begin{exercice}[3]{Choisir le côté d’une dalle}
Une dalle carrée a une aire de 2 m$^2$. Donner son côté exact en mètres, puis l’arrondir au millimètre. Peut-on le remplacer par 1,4 m pour calculer une aire exacte ?
\end{exercice}\end{blocLivre}
\begin{blocLivre}{ex-2-15}\begin{exercice}[3]{Un cercle et trois droites}
Dans un repère orthonormé, le cercle de centre $O$ et de rayon 5 contient les points $M(x\,;y)$ tels que $OM=5$, c’est-à-dire $x^2+y^2=25$ d’après Pythagore. On le coupe par les droites horizontales d’équations $y=3$, $y=5$ et $y=6$. Pour chacune, résoudre l’équation obtenue en $x$ et conclure.
\label{fig-ex-2-15}\figureGeo{cercle-droites}
\end{exercice}\end{blocLivre}

\clearpage
